\appendix
\section{Proof of \protect{\expref{Theorem}{basic theorem}}}

The goal of the current section is the prove 
\expref{Theorem}{basic theorem}.
We remind the reader that this result
follows from the characterization by
Kasami and Tokura~\cite{kt70} of 
all codewords of the Reed-Muller code
that have weight at most twice the
minimal weight.

First note that the bound of \expref{Theorem}{basic theorem} is
sharp as it is obtained by
$$x^\alpha (x^\beta+x^\gamma)$$ where $\alpha$,
$\beta$, and $\gamma$ are disjoint multi-indices
where the size of $\alpha$ is $r$ and the sizes
of $\beta$ and $\gamma$ both are $d-r$.

\begin{proof}
We prove the statement by induction and
we establish $d=2$ as the base case below
to avoid degenerate cases.  The following
easy observations are useful for us.

\begin{itemize}
\item As $d-r$ cannot equal 1, it follows
that for any degree-$d$ polynomial that
is not a product of affine factors, the
fraction of inputs on which it takes
the value one is at least $(3/2)\, 2^{-d}$.

\item The lower bound to be proved never exceeds 
$2^{1-d}$, and thus this bound is
always sufficient (but not always possible).
\end{itemize}
The statement for general $d$ and $r$ follows
from the case of $d-r$ and $0$ and hence we
may focus on the case of $r=0$ (but of course
use arbitrary $r$ in the inductive statements).

A fact that needs to be taken into account is, as we
are using that $x_i^2=x_i$,
that the factorization is not unique.  In
particular if $A$ and $A'$ are
two affine factors of a polynomial $P$ then another
factor is $1+A+A'$
as 
$$
(1+A+A')A=A'A\,.
$$
Thus if we have several affine factors
we can construct new affine factors by taking
the sum of of an even number of such factors
added with the constant 1 or the sum of
an odd number of such factors. 
Let us point out that the statement
of \expref{Theorem}{basic theorem} that $Q$ does not
contain any affine factor is intended to mean that there is no
factorization of $P$ that contains the given
affine factors and one more linearly independent
affine factor.

Another useful
fact is that an affine function $A(x)$ is a factor of 
a polynomial $P$
iff $A(x)=0$ implies $P(x)=0$ or equivalently
if $P(y)=1$ implies $A(y)=1$.  The non-obvious
direction of this statement was established during
the proof of \expref{Theorem}{thm upper sat}
when identifying implied affine constraints.

For the reader worried about the non-unique factorization,
let us point out the following fact that we leave to
the reader to verify.  The number of affine factors
is the co-dimension of the affine hull of all
points such that $P(x)=1$.  The factors that
may appear in the factorization are the affine equations
defining this affine hull.  In the full factorization
we may take any full set of linearly independent
equations.

Let us address the case of $d=2$, which can
be established by the normal form of 
%
degree-two 
polynomials, but let us follow a
different path to prepare for the general proof.
As stated above, the interesting case is
when $P$ has no affine factor and let
us write $P(x)=P_0(x)+x_1P_1(x)$.

Consider setting $x_1$ to its two values and as
$P$ does not contain an affine factor neither
of the induced polynomials can be identically 0.  Furthermore
if neither of these settings result in a polynomial
with an affine factor we are done by
induction.  Let us finally assume that $x_1=0$
results in an affine factor, which we, by an
affine change of coordinates, can assume is $x_2$.
Thus we can assume that 
$$
P(x)= x_2 A_2(x)+x_1A_1(x)\,,
$$
for two affine functions $A_1$ and $A_2$ where
we also can assume that $x_i$ does not appear in $A_i(x)$ as it
can be replaced by 1 giving the same result.
The main case is that the collection of $x_1,x_2, A_1(x),$
and $A_2(x)$ form linearly independent affine functions and
in this case the fraction of inputs for which
$P$ is one is exactly $3/8$ which is the claimed bound.
We have a number of cases to consider when the four
functions are not linearly independent.

\begin{enumerate}
\item $A_1(x) \equiv 1$.  If $A_2(x)=x_1$, then $x_1$ is a factor
of $P$ while if $A_2(x)=1+x_1$ then $P$ is one with
probability $3/4$. Finally if $A_2(x)$ is linearly independent
of $x_1$ this probability is $1/2$.  

\item $A_1(x)=x_2$ makes $x_2$ a factor of $P$.

\item $A_1(x)=1+x_2$ makes $Pr[P(x)=1]=1/2$ unless we
have $A_2(x) \equiv 1$ when this probability is $3/4$.

\item $A_1(x)$ is linearly independent of $x_2$.  In this case,
by an affine change of variables, we can assume that $A_1(x)=x_3$.
Now since $A_2(x)$ does not contain $x_2$, is linearly
dependent of $x_1$ and $x_3$, and is not a factor
of $x_1x_3$ (ruling out also $(1+x_1+x_3)$), 
$A_2(x)$ must equal one of the functions 1, $1+x_1$
or $1+x_3$.  In each of these cases it is easy
to check that $Pr[P(x)=1]$ is $1/2$.

\end{enumerate}
This finishes the case $d=2$ end we turn to the
general case.  Not surprisingly, also here we
end up analyzing a number of cases.

As in the case $d=2$ we can assume that $P$
has no affine factors but the polynomial resulting
when substituting $x_1=0$ gives a polynomial
with at least one affine factor.  Picking
a full set of linearly independent factors
and making an affine transformation we can
assume that
$$
P(x)=x_2 x^\beta P_2(x)+ x_1 P_1(x)$$
where $\beta$ is a  possible empty multi-index and
$P_2$ has no affine factors. First
let us consider affine factors in $P_1$.  

Let $\prod_{i=1}^r A_i(x)$ be the affine
factors that appear in $P_1$.  We have two
cases depending whether each $A_i$ that might appear
in the factorization, together with the
affine forms that appear in the first product
(i.\,e., the coordinate functions given by $x_2$
and the elements of $\beta$) are linearly independent.

Suppose these functions are not linearly independent and hence that
we can choose the factorization such
that $A_1(x)$ only depends
on $x_2$ and the variables in $\beta$.  Let
us look at the point $x^0$ where $x_2$ and all
elements of $\beta$ equals $1$.  If $A_1(x^0)=1$
then $A_1(x)$ is a factor of $P$
and this is a contradiction of the assumptions.
If, on the other hand $A_1(x_0)=0$ then the sets of points
where $x_2 x^\beta P_2(x)=1$ and $x_1 P_1(x)=1$
are disjoint and, as each of these sets is of density
at least $2^{-d}$,  we get $\Pr[P(x)=1] \geq 2^{1-d}$ and the
theorem follows also in this case.

Now assume that any affine form  appearing in the factorization
of $P_1$ is linearly independent of $x_2$ and the variables
in $\beta$.  Then, by an affine transformation,
we can assume that 
\begin{eqnarray}
P(x)=x_2 x^\beta P_2(x)+ x_1 x^\gamma P_1'(x)\,,
\end{eqnarray}
where $\beta$ and $\gamma$ are disjoint
multi-indices and no more affine factors
can be pulled out of $P_2$ or $P_1'$.

Let us analyze what happens for
the four possible simultaneous assignments
of values to $x_1$ and $x_2$.
When both are 0 we get a function that is
identically 0 which is not good for us
but in the other cases we get non-zero polynomials
of degree $d-1$ and we now analyze the 
structure of these polynomials.

When $x_1=0$ and $x_2=1$
we get $x^\beta P_2(x)$, when
$x_1=1$ and $x_2=0$ we get
$x^\gamma P_1'(x)$, and in the final
case we have
$$
W(x)=x^\beta P_2(x)+ x^\gamma P_1'(x)\,.$$
Note that $W$ is not identically 0
as $(x_1+x_2)$ then would have been
an affine factor of $P$ and furthermore that
$x^{\beta}P_2$, $x^{\gamma}P_1'$, $W$ all are of degree at
most $d-1$.

Suppose first that neither $P_2$ nor $P_1'$ is
the constant one.
Then, using the bound that a 
%
degree-$(d-1)$
polynomial that
is not a pure product of affine factors is one with
probability at least $(3/2)\,2^{1-d}$, we have
that $\Pr[P(x)=1]$ is at least
$$
\frac 14 \left(\frac 32 2^{1-d}+\frac 32 2^{1-d}+2^{1-d}\right)= 2^{1-d}
$$
proving the bound in this case.
By symmetry we may thus assume 
that $P_2\equiv 1$.
If $\gamma=\emptyset$
or $W$ does not contain any affine factor,
then $\Pr[P(x)=1]$ is at least
$$
\frac 14 (2^{1-d}+2^{1-d}+2^{2-d}-2^{3-2d})\,,
$$
which exactly equals the claimed bound.
Thus we can assume that $\gamma$ is non-empty
and $W$ contains an affine factor $A(x)$ and remember
that \begin{eqnarray}\label{weq}
W(x)=x^\beta + x^\gamma P_1'(x)\,.
\end{eqnarray}

Suppose $A$ only depends
on variables in $\beta$.  If it is fixed to
1 by setting all variables in $\beta$
to one, then $A(x)$ is a factor of $x^\beta$
and hence also of $W(x)+x^\beta=x^\gamma P_1'(x)$
and hence also of $P(x)$, contradicting assumptions.
If $A(x)$ is forced to 0 by this assignment
then setting any variable
in $\gamma$ (remember it is non-empty and disjoint
from $\beta$) to 0 and we get $W(x)=0$ while
$x^\beta=1$ and $x^\gamma P_1'(x)=0$ contradicting
\eqref{weq}.  Thus we can assume
that $A(x)$ depends on some variable
outside $\beta$.

If we can fix some variable in $\gamma$ to 0,
the variables of $\beta$ to one and $A$ to zero
we get a contradiction to \eqref{weq} as
$x^\beta=1$ while $W(x)=0$ and $x^\gamma=0$.
If the size of $\gamma$ is
at least 2 or $W$ contains at least two different
affine factors we claim that this must be possible.  Suppose first
that the size of $\gamma$ is at least 2.

As $A(x)$ does not only depend on variables
in $\beta$ we can fix these variables to
one, and then pick a suitable variable
in $\gamma$ to fix to 0 without fixing
the value of $A(x)$. We can then fix
additional variables to make $A(x)=0$
obtaining the desired contradiction.

Now suppose that $W$ contains at least
2 affine factors and let $A'$ be a factor distinct
from $A$.  In this case,
one of $A$, $A'$ and $1+A+A'$ is
a factor of $W$ and does not
depend on the first variable of $\gamma$
(and some variable outside $\beta$).
It follows again that we can
make $x^\beta=1$, $x^\gamma=0$ and $W(x)=0$,
again obtaining a contradiction.

The only remaining case is when $\gamma$ is
of size one and $W$ has one affine factor.
In this case, by induction
$\Pr[P(x)=1]$ is at least
$$
\frac 14 \left(2^{1-d}+ 2\cdot \frac 12 \left(2^{3-d}-2^{5-2d}\right)\right).
$$
For $d$ at least 4 this is at least $2^{1-d}$
and we are done.  Finally note that in the
case $d=3$, $P_1'$ as well as the co-factor
of $A$ in $W$ are of degree at most one
and hence if they do not contain an
additional factor they must be the constant
1 and the theorem follows also in this case.
\end{proof}
